% SIMRUN Project Documentation
\documentclass[11pt,a4paper]{article}
\usepackage[utf8]{inputenc}
\usepackage[T1]{fontenc}
\usepackage{lmodern}
\usepackage{geometry}
\usepackage{hyperref}
\usepackage{longtable}
\usepackage{tabularx}
\usepackage{graphicx}
\usepackage{float}
\usepackage{verbatim}
\usepackage{amsmath}
\geometry{margin=1in}

% Document metadata
\title{SIMRUN -- Project Documentation\\\large Simulation of Really Unpredictable Nonsense}
\author{SIMRUN Contributors}
\date{\today}

\begin{document}

% Title page
\maketitle
\thispagestyle{empty}
\vfill
\begin{center}
  \textbf{Repository overview}\\
  Top-level layout: `src/` contains `ui/`, `compiler/`, `sim/`, `analysis/` and tooling such as `vcpkg/`.
\end{center}
\newpage

% Table of contents (index)
\tableofcontents
\newpage

% Start of main documentation
\section{Problem Statement}

SIMRUN is a research and engineering project whose goal is to provide a flexible environment to design, compile and run architecture-driven simulations of distributed systems. The project name expands as: \emph{Simulation of Really Unpredictable Nonsense}.

The core problems addressed by SIMRUN are:
- Provide a visual UI to design component-based architectures and simulation scenarios.
- Serialize architecture designs into a stable intermediate representation (IR) consumable by a C++ compiler/backend.
- Offer a reproducible simulation engine that executes compiled scenarios and returns structured results.
- Enable rapid development and iterative testing via an Electron-wrapped UI with hot-reload.

\section{Deliverables}

This documentation describes the deliverables that accompany the repository and expected artifacts after a full build:

- UI: A React + TypeScript application, packaged and runnable via Electron. Development mode supports hot-reload and live simulation calls.
- Compiler/Server: A C++ component (built with CMake and vcpkg) which accepts serialized IR, compiles or interprets it and runs simulation tasks. It may expose a local server interface used by the UI to request simulations.
- Simulation Engine: C++ modules implementing domain entities, factories and the runtime used to execute the generated scenarios.
- Analysis Tools: Utility programs that inspect `final_ir.json` and other intermediate results for verification, metrics and visualization.
- Documentation: This LaTeX document and source README/HOWTO files providing build and run instructions.

\newpage

\section{Repository Layout and Responsibilities}

The repository is organized to separate concerns between frontend, compiler/backend and simulation runtime. The high-level layout:

\begin{itemize}
  \item `src/ui/` -- React + TypeScript UI, Vite build system, Electron wrapper, UI components and stores.
  \item `src/compiler/` -- C++ compiler backend; CMake project. Consumes serialized IR and produces runnable artifacts or orchestrates simulations.
  \item `src/sim/` -- C++ simulation engine; contains `core/`, `entities/`, `events/`, `factory/` and other runtime modules.
  \item `src/analysis/` -- Analysis and tooling to inspect IR and final outputs.
  \item `vcpkg/` -- vcpkg bootstrap and portfiles used for C++ dependency management.
\end{itemize}

\subsection{Key files referenced}
\begin{longtable}{p{0.25\linewidth} p{0.7\linewidth}}
`README.md` & Project title and short overview.\\
`HOWTO_RUN.txt` & Step-by-step developer guide for UI-only and full (UI + compiler) workflows.\\
`src/ui/package.json` & Shows UI scripts such as `electron:dev` and `electron`.\\
`src/ui/src/pages/Index.tsx` & Main page that composes the canvas and runs simulation by sending IR to the compiler.\\
`src/compiler/CMakeLists.txt` & Top-level CMake configuration for the compiler backend.\\
`src/sim/` & C++ simulation entry points and modules (engine, entities, factories).\\
\end{longtable}

\newpage

\section{UI Frontend}

The UI is implemented using React + TypeScript with Vite for fast iteration. Key technologies include Tailwind CSS, shadcn-ui components, and a state management layer exposed as stores.

\subsection{Main UX flow}

Users design architectures on a canvas using drag-drop components. Each component includes a `category` (e.g., `api`, `database`, `cache`, `network`) and parameter presets. The UI composes an export object containing components, links, routes, workload and faults.

\subsection{Example: serialization and run path}
The UI performs the following steps when a user clicks \textit{Run Simulation} (sequence observed in `src/ui/src/pages/Index.tsx`):

\begin{enumerate}
  \item Gather components (`nodes`) and links (`edges`) from the architecture store.
  \item Compose an export object with components, links, routes, workload and metadata.
  \item Call `serializeToIR(exportData)` to generate an intermediate representation.
  \item Wrap the IR into a compiler request with `wrapForCompiler(ir)`.
  \item Invoke `window.api.simulate(compilerRequest)` using Electron's IPC/native bridge.
\end{enumerate}

\subsection{Scripts and development commands}
The UI `package.json` includes the following helpful npm scripts:
\begin{verbatim}
  "dev": "vite",
  "build": "vite build",
  "build:dev": "vite build --mode development",
  "lint": "eslint .",
  "preview": "vite preview",
  "electron": "electron .",
  "electron:dev": "concurrently \"npm run dev\" \"wait-on http://localhost:5173 && electron .\""
\end{verbatim}

`electron:dev` starts the Vite dev server and launches Electron once the UI host is ready, enabling hot reload.

\newpage

\section{Electron Integration}

Electron is used to package the web UI as a desktop application and to provide a secure IPC layer to the C++ components. The UI uses an exposed API on `window.api` for requesting simulations. Typical flow:

\begin{itemize}
  \item Development: `npm run electron:dev` (starts Vite and Electron together).
  \item Production: `npm run build` then `npm run electron` packages and launches the app.
\end{itemize}

Electron bridges the renderer (React) and native capabilities (filesystem, child processes) and is therefore the place to route `simulate` requests from the renderer to either a bundled compiler binary or a local compiler server.

\section{Compiler / Server (C++)}

The compiler backend is a C++ CMake project under `src/compiler/`. It is responsible for:

\begin{itemize}
  \item Validating and parsing the IR produced by the UI.
  \item Translating IR to simulation instructions or native code consumed by the simulation engine.
  \item Optionally exposing an HTTP/IPC server to accept simulation requests (the HOWTO references a port such as 8081).
  \item Returning structured results to the UI (JSON-like payloads) containing event traces, metrics and error reports.
\end{itemize}

Build notes (excerpted from HOWTO_RUN.txt): use CMake and vcpkg toolchain; on Windows prefer the Visual Studio generator; on Linux set `CMAKE_BUILD_TYPE=Release` and build with `-j$(nproc)`.

\subsection{Server considerations}
If the compiler exposes a long-lived server (for quicker iteration), ensure:

\begin{itemize}
  \item Authentication is restricted to local requests only (bind to localhost).
  \item Service detects and reports malformed IR with helpful diagnostics.
  \item Process isolation when invoking heavy simulation tasks to avoid blocking the server thread.
\end{itemize}

\newpage

\section{Simulation Engine: `src/sim/`}

The simulation engine contains the runtime implementation of the component models, event scheduling and metrics collection. Typical subdirectories include:

\begin{itemize}
  \item `core/` -- central scheduler, timebase, common utilities.
  \item `entities/` -- component implementations (API servers, databases, caches, network links).
  \item `events/` -- event types and handlers used by the scheduler.
  \item `factory/` -- factories responsible for constructing entity instances based on IR entries.
\end{itemize}

The factory mechanism maps the IR `components` entries into runtime entity instances with parameterized behavior. This supports swapping implementations (for example, different database models) while reusing the same IR.

\subsection{Entity design notes}

Entities should follow these design principles:

\begin{itemize}
  \item Clear separation between configuration (parameters) and runtime state.
  \item Well-defined event interfaces for communication (messages, requests, responses).
  \item Deterministic scheduling when required for reproducible experiments.
\end{itemize}

\newpage

\section{IR (Intermediate Representation) and APIs}

The IR is the canonical interchange format between the UI and the compiler/simulation backend. It typically contains:

\begin{itemize}
  \item `context.components`: list of component descriptors (id, type, parameters, profile).
  \item `links`: network connections and their parameters.
  \item `routes`, `workload`, `faults`: scenario-level specifications.
  \item `metadata`: version, timestamp and provenance information.
\end{itemize}

Example IR flow from `Index.tsx`:
\begin{verbatim}
  const ir = serializeToIR(exportData);
  const compilerRequest = wrapForCompiler(ir);
  window.api.simulate(compilerRequest)
\end{verbatim}

When designing the IR, prefer a stable schema (use a versioned `metadata.version`) and include explicit validation rules to keep evolution manageable.

\section{Build and Run Instructions (Concise)}

The `HOWTO_RUN.txt` file documents both UI-only and full setups. Minimal quick start (UI only):

\begin{verbatim}
cd src/ui
npm install
npm run electron:dev
\end{verbatim}

Full build (compiler + UI) generally follows the steps below (Windows flavor):

\begin{verbatim}
cd src/compiler
mkdir build && cd build
cmake -G "Visual Studio 17 2022" -A x64 -DCMAKE_TOOLCHAIN_FILE=..\..\vcpkg\scripts\buildsystems\vcpkg.cmake ..
cmake --build . --config Release
\end{verbatim}

After building the compiler, run the UI as above and use the Electron app to trigger simulations.

\newpage

\section{Appendix: Design Recommendations and Next Steps}

This project is organized to support iterative development. Recommended next steps and best practices:

\begin{itemize}
  \item Add schema tests for the IR and unit tests for `serializeToIR` and `wrapForCompiler`.
  \item Add CI steps to build the compiler with vcpkg and to run UI linting and type checks.
  \item Document the compiler server API (endpoints, request/response examples and error codes).
  \item Provide example scenario files (IR samples) under `docs/examples/` for quick experimentation.
\end{itemize}

\vfill
\begin{center}
  Generated documentation for SIMRUN. For more details see the repository `README.md` and `HOWTO_RUN.txt`.
\end{center}

\end{document}
